前面的 Pauli--Lubanski 算符已经告诉我们，质量和自旋可以作为单粒子态的 Lorentz 不变量标签。但还缺少一个更具体的问题：把粒子从一个惯性系变到另一个惯性系时，四动量会改变，内部自旋指标究竟怎样跟着变？

固定一条质量壳后，任意四动量 $p$ 都可由某个标准 Lorentz 变换从参考四动量 $\bar p$ 得到。若再施加一般 Lorentz 变换 $\Lambda$，把 $\Lambda p$ 重新拉回 $\bar p$ 所留下的复合变换会保持 $\bar p$ 不变，因此它不能改变参考动量，只能作用在内部指标上。所有保持 $\bar p$ 不变的 Lorentz 变换组成\emph{小群}。有质量粒子的自旋和无质量粒子的螺旋度，正是这个剩余对称群的不可约表示标签。

为使 Lorentz 变换与单粒子完备关系兼容，单粒子动量本征态采用协变归一化。对内部量子数 $r$，取
\begin{equation}
 \boxed{
 \langle \mathbf p',r'|\mathbf p,r\rangle_{\rm cov}
 =2p^0(2\pi\hbar)^3\delta_{r'r}
 \delta^{(3)}(\mathbf p'-\mathbf p),
 \qquad p^0=\frac{E_{\mathbf p}}{c},}
 \label{eq:rel-state-normalization-si-r9}
\end{equation}
并配套使用
\begin{equation}
 \boxed{
 \id(1)=\sum_r\int
 \frac{\dd^3 p}{(2\pi\hbar)^3\,2p^0}
 |\mathbf p,r\rangle_{\rm cov}
 {}_{\rm cov}\langle\mathbf p,r|.}
 \label{eq:rel-completeness-si-r9}
\end{equation}
这组归一化使测度 $\dd^3 p/(2p^0)$ 保持 Lorentz 不变，因此 Wigner 旋转的内部矩阵可以与动量测度分离，不会额外出现由非协变三维 $\delta$ 归一化产生的 Jacobian 平方根。

固定一条给定质量壳，并在这条轨道上选取标准动量 $\bar p^\mu$。对每个同轨道四动量 $p^\mu$，选一个标准 Lorentz 变换 $L(p)$ 满足
\begin{equation}
 L(p)\bar p=p.
 \label{eq:wigner-standard-boost-dp9}
\end{equation}
以 $|\bar p,\sigma\rangle$ 表示标准动量上的的内部基，定义
\begin{equation}
 |p,\sigma\rangle
 \equiv U[L(p)]|\bar p,\sigma\rangle.
 \label{eq:wigner-state-definition-dp9}
\end{equation}
$\sigma$ 标记固定标准动量时仍可变化的内部自由度，不增加时空坐标。对任意 Lorentz 变换 $\Lambda$，定义
\begin{equation}
 \boxed{
 W(\Lambda,p)
 \equiv L^{-1}(\Lambda p)\,\Lambda\,L(p),
 \qquad
 W(\Lambda,p)\bar p=\bar p.}
 \label{eq:wigner-little-group-element-dp68}
\end{equation}
因此 $W(\Lambda,p)$ 不再改变四动量，只能作用在内部指标上。若内部空间采用小群不可约表示 $D$，单粒子态的完整 Lorentz 变换为
\begin{equation}
 \boxed{
 U(\Lambda)|p,\sigma\rangle
 =\sum_{\sigma'}
 D_{\sigma'\sigma}[W(\Lambda,p)]
 |\Lambda p,\sigma'\rangle.}
 \label{eq:wigner-state-transform-dp68}
\end{equation}

\begin{figure}[htbp]
 \centering
 \begin{tikzfig}[x=1cm,y=1cm,>=Stealth]
\node[lecturebox] (k) at (-4.5,0) {standard momentum\\$\bar p$};
\node[lecturebox] (p) at (-1.2,2.0) {$p=L(p)\bar p$};
\node[lecturebox] (lp) at (3.1,2.0) {$\Lambda p$};
\node[lecturebox] (kb) at (3.1,-1.0) {$\bar p$\\same orbit reference};

\draw[lecturearrow,wfA] (k)--node[line-label,above left] {$L(p)$}(p);
\draw[lecturearrow,wfA] (p)--node[line-label,above] {$\Lambda$}(lp);
\draw[lecturearrow,wfA] (lp)--node[line-label,right] {$L^{-1}(\Lambda p)$}(kb);
\draw[lecturearrow,wfA,dashed] (k)--node[line-label,below] {$W(\Lambda,p)$}(kb);

\node[lecturebox,align=center,font=\normalsize] (stab) at (-0.7,-2.65)
  {$G_{\bar p}=\{\Lambda\mid \Lambda\bar p=\bar p\}$\\stabilizer of the reference momentum};
\node[lecturebox] (massive) at (-3.2,-4.65) {$\bar p^2=m^2c^2>0$\\$G_{\bar p}\simeq SO(3)$\\double cover $SU(2)$\\spin $s$};
\node[lecturebox] (massless) at (3.4,-4.65) {$\bar k^2=0$\\$G_{\bar k}\simeq ISO(2)$\\discrete-helicity sector\\helicity $\lambda$};
\draw[lecturearrow,wfA] (stab.south west) --node[line-label,left,font=\normalsize]{timelike orbit}(massive.north);
\draw[lecturearrow,wfA] (stab.south east) --node[line-label,right,font=\normalsize]{null orbit}(massless.north);
\end{tikzfig}

 \caption{Wigner 小群构造。$L(p)$ 把标准动量推到 $p$，而 $W(\Lambda,p)=L^{-1}(\Lambda p)\Lambda L(p)$ 把复合变换拉回同一个标准动量，因此只作用在内部自旋/螺旋度空间。}
 \label{fig:wigner-little-group-dp9}
\end{figure}

对 $m>0$，最方便的标准动量就是静止四动量 $\bar p_m^\mu=(mc,\mathbf0)$。保持它不变的 Lorentz 变换不能混合时间与空间分量，只能在三维空间做转动，因此
\begin{equation}
 \boxed{
 G_{\bar p_m}\simeq SO(3),
 \qquad
 \widetilde G_{\bar p_m}\simeq SU(2).}
 \label{eq:wigner-massive-little-group-dp68}
\end{equation}
量子态使用双覆盖 $SU(2)$，所以其不可约表示由 $s=0,\tfrac12,1,\ldots$ 标记，并具有 $2s+1$ 个自旋态。这与前文 $W^2=-m^2c^2\hbar^2s(s+1)$ 的 Casimir 分类是同一件事的两种表述。

无质量轨道没有静止系，因此不能把有质量构造简单令 $m\to0$。选沿 $+z$ 传播的标准零质量动量
\begin{equation}
 \bar k^\mu=(\kappa_\star,0,0,\kappa_\star),
 \qquad \kappa_\star>0,
 \qquad \bar k^2=0,
 \label{eq:wigner-massless-standard-momentum-dp9}
\end{equation}
其中 $\kappa_\star$ 具有物理动量量纲，其数值只是在零质量轨道上选参考点。全部稳定子方向由满足 $X\bar k=0$ 的 Lorentz Lie 代数元给出。写一般元
\begin{equation*}
 X=\sum_{i=1}^3(a_iJ_i+b_iK_i)
\end{equation*}
出发，并要求它在 $P^\mu=\bar k^\mu$ 上满足 $[X,P^\mu]=0$。四个分量的线性条件化为
\begin{equation*}
 b_3=0,
 \qquad
 b_1=-a_2,
 \qquad
 b_2=a_1,
\end{equation*}
所以稳定子 Lie 代数恰为三维。取一组方便的基为
\begin{equation}
 \mathsf N_1\equiv K_1-J_2,
 \qquad
 \mathsf N_2\equiv K_2+J_1,
 \qquad
 J_3.
 \label{eq:wigner-massless-little-group-generators-dp68}
\end{equation}
它们满足
\begin{equation}
 \boxed{
 [J_3,\mathsf N_1]=\ii\hbar\mathsf N_2,
 \qquad
 [J_3,\mathsf N_2]=-\ii\hbar\mathsf N_1,
 \qquad
 [\mathsf N_1,\mathsf N_2]=0,}
 \label{eq:wigner-massless-little-group-algebra-dp68}
\end{equation}
这正是二维平面转动加平移的代数，因此在 $SO^+(1,3)$ 中零质量标准动量的小群同构于 $ISO(2)$。这里的“平移”作用在内部表示空间，不是把粒子的时空位置平移。量子态若由 Lorentz 群的双覆盖 $SL(2,\mathbb C)$ 实现，则实际作用的是 $ISO(2)$ 的相应双覆盖；这一区分对半整数螺旋度是必要的。

这个小群的量子不可约表示有两类。若 $\mathsf N_{1,2}$ 非平凡作用，会得到连续自旋表示；若这两个“平移型”生成元在物理态上平凡作用，则只剩绕运动方向的转动生成元 $J_3$。这种第二类表示由一个离散或半整数的螺旋度标签 $\lambda$ 区分，并满足
\begin{equation}
 \mathsf N_1|\bar k,\lambda\rangle=0,
 \qquad
 \mathsf N_2|\bar k,\lambda\rangle=0,
 \qquad
 J_3|\bar k,\lambda\rangle=\hbar\lambda|\bar k,\lambda\rangle.
 \label{eq:massless-helicity-irrep-dp9}
\end{equation}
此时 Pauli--Lubanski 矢量与四动量平行：
\begin{equation}
 \boxed{W^\mu|k,\lambda\rangle
 =\hbar\lambda\,P^\mu|k,\lambda\rangle.}
 \label{eq:massless-pauli-lubanski-helicity-dp9}
\end{equation}
因此无质量表示的 $W^2=0$ 对所有离散螺旋度都相同；真正区分内部态的是小群转动生成元本征值 $\lambda$。光子的物理 Hilbert 空间取 $\lambda=\pm1$，无质量自旋 $1/2$ 场取 $\lambda=\pm\tfrac12$。光子的 $\lambda=0$ 不属于物理横向单光子空间；后面 Maxwell 约束与规范商空间会从场变量层面再次得到同一结论。

三个常被混写的对象属于不同的表示层级：$(j_L,j_R)$ 是\emph{局域场}在有限维 Lorentz 复表示中的变换类型；$s$ 是有质量单粒子小群双覆盖 $SU(2)$ 的不可约标签；$\lambda$ 是无质量离散表示中稳定子转动方向的本征值。对整数自旋可直接把它看成 $SO(2)$ 角色标签；对半整数自旋必须在 Lorentz 双覆盖中的小群上理解它。一个四矢量场属于 $(\tfrac12,\tfrac12)$ 并不意味着物理光子有四个自旋态；规范约束最终只留下 $\lambda=\pm1$ 两个单粒子螺旋度态。

